\documentclass{beamer}

\usepackage{cuhkbeamer}

\begin{document}

% ====================== Slide 1: TITLE SLIDE (REVISED) ======================
\begin{frame}
  \title{LAE Policy and Regulation}
  \author{Chiu Yu KO}
  \institute{Low Altitude Economics}
  \date{Week 7}
  \titlepage
\end{frame}


% ====================== Slide 2: FORMAT & OBJECTIVES ======================
\begin{frame}[t]
\frametitle{Learning Objectives}
By the end of this class, you should be able to:
\begin{itemize}
\item \textbf{Explain} Hong Kong's risk-based framework for small unmanned aircraft and the role of Sandbox X.
\item \textbf{Identify} the evidence and permissions required for advanced operations, including BVLOS testing.
\item \textbf{Evaluate} how safety assurance affects approval, liability, reliability, and insurability.
\item \textbf{Explain} how digital traffic-management services may support safe and scalable operations.
\item \textbf{Design} a phased regulatory roadmap that links operating privileges to evidence and performance.
\end{itemize}
\end{frame}

% ====================== Slide 3: ICEBREAKER ======================
\begin{frame}[t]
\frametitle{The Regulatory Strategy Meeting}
\textbf{Imagine an operator preparing to test a complex drone route in Hong Kong.}
\begin{itemize}
\item The operating team must explain the aircraft, route, people, procedures, communications, contingency measures, and data to be collected.
\item The regulator must decide whether the proposed controls are proportionate to the operational risk.
\end{itemize}
\vspace{0.3cm}
Regulation can increase cost and delay, but it can also establish safety, trust, predictable market access, and credible evidence.
\begin{itemize}
\item \textbf{Quick poll (2 min):} Which is more valuable to an early operator: faster approval, clearer requirements, or permission to test a more complex operation?
\end{itemize}
\end{frame}

% ====================== Session 1 Title Card ======================
\begin{frame}[c]
  \begin{center}
    \huge\color{cuhkpurple}\textbf{Session 1}\\
    \vspace{0.5cm}
    \Large Regulatory Moats, Sandbox X \& BVLOS
  \end{center}
\end{frame}

% ====================== Slide 4: CONCEPT PRIMER 1 - ECONOMIC MOAT ======================
\begin{frame}[t]
\frametitle{Concept Primer: Regulation and Competitive Advantage}
Regulatory capability may become a competitive advantage when a firm can demonstrate safe, reliable, and compliant operations more effectively than its rivals.
\begin{itemize}
\item Approval can require investment in people, procedures, engineering, testing, documentation, and monitoring.
\item Evidence and operating experience may reduce uncertainty and improve credibility with customers, partners, investors, and insurers.
\item However, regulation does not grant an automatic monopoly, and approvals should not be treated as tools for excluding efficient competitors.
\item Requirements may also promote entry by creating clear, transparent, and proportionate pathways.
\end{itemize}
\textbf{Managerial lesson:} Build a compliance capability, not a strategy based on presumed regulatory protection.
\end{frame}

% ====================== Slide 5: APPLIED TOOL - ECONOMIC MOAT VALUE ======================
\begin{frame}[t]
\frametitle{Applied Tool: Value of Earlier Market Readiness}
\begin{block}{A Simple Screening Relationship}
\[
\text{Net Readiness Value}
=
\text{Contribution Earned During the Lead Period}
-\text{Additional Compliance and Launch Cost}
\]
\end{block}
\begin{itemize}
\item The lead period begins only when revenue-generating operation is actually permitted and demand exists.
\item Contribution, not revenue or accounting profit, should reflect additional operating costs.
\item Faster readiness may create learning, reputation, and partnership benefits, but competitors and substitutes can still enter.
\item The calculation should include probability of approval, delay, unsuccessful trials, and the possible expiry of the lead.
\end{itemize}
\end{frame}

% ====================== Slide 6: BACK-OF-THE-ENVELOPE - ECONOMIC MOAT ======================
\begin{frame}[t,shrink=10]
\frametitle{Back-of-the-Envelope: The Value of Earlier Readiness}
\textbf{Illustrative assumptions:} Monthly contribution after approval = HK\$2M; additional compliance cost = HK\$15M.
\begin{columns}[T]
\begin{column}{0.48\textwidth}\begin{block}{Twelve-Month Lead}
\begin{itemize}
\item Contribution during lead: $12\times2=\text{HK\$24M}$
\item Additional compliance cost: HK\$15M
\item \textbf{Net readiness value: HK\$9M}
\end{itemize}\end{block}\end{column}
\begin{column}{0.48\textwidth}\begin{block}{Six-Month Lead}
\begin{itemize}
\item Contribution during lead: $6\times2=\text{HK\$12M}$
\item Additional compliance cost: HK\$15M
\item \textbf{Net readiness value: $-$HK\$3M}
\end{itemize}\end{block}\end{column}
\end{columns}
\textit{Takeaway: Early readiness has value only if the lead is sufficiently long and commercially useful. It does not imply monopoly protection.}
\vfill{\tiny Illustrative assumptions; excludes probability of approval, financing, tax, and demand uncertainty.}
\end{frame}

\begin{frame}[t]
\frametitle{Concept Primer: Hong Kong's Regulatory Sandbox X}
Sandbox X enables selected project concepts to be tested on predefined routes and under controlled conditions.
\begin{itemize}
\item It supports testing of airspace management, flight operations, infrastructure, communications, and emergency response.
\item It allows government and participants to collect evidence on safety, reliability, resilience, technical feasibility, and commercial viability.
\item Cross-boundary trials may also examine coordination and customs, immigration, and quarantine procedures.
\item Projects proceed progressively, with routes, parameters, and timing subject to approved conditions.
\end{itemize}
\textit{Sandbox X is not a legal safe space, a waiver of safety obligations, or automatic authorization for commercial service.}
\end{frame}
% ====================== Slide 8: APPLIED TOOL - SANDBOX X VALUE ======================
\begin{frame}[t]
\frametitle{Applied Tool: Value of a Controlled Trial}
\begin{block}{A Simple Expected-Learning Calculation}
\[
\text{Trial Value}
=
\text{Expected Avoided Redesign and Delay Cost}
+\text{Value of Better Information}
-\text{Trial Cost}
\]
\end{block}
\begin{itemize}
\item A trial can reveal whether routes, communications, infrastructure, procedures, and customer demand work as expected.
\item Evidence may support later approval or investment, but the timing and outcome are not guaranteed.
\item Trial flights and commercial operations are different stages and may involve different permissions and conditions.
\item The principal benefit is earlier, lower-cost learning rather than presumed early monopoly profit.
\end{itemize}
\end{frame}

% ====================== Slide 9: BACK-OF-THE-ENVELOPE - SANDBOX X ======================
\begin{frame}[t,shrink=10]
\frametitle{Back-of-the-Envelope: The Value of Earlier Learning}
\textbf{Illustrative assumptions:} A controlled trial costs HK\$5M. Without a trial, there is a 40\% chance that a HK\$20M redesign is discovered after investment.
\begin{columns}[T]
\begin{column}{0.48\textwidth}\begin{block}{No Early Trial}
\begin{itemize}
\item Expected redesign cost: $0.40\times20$
\item \textbf{HK\$8M}
\item Information arrives after larger commitment.
\end{itemize}\end{block}\end{column}
\begin{column}{0.48\textwidth}\begin{block}{Controlled Trial}
\begin{itemize}
\item Trial cost: HK\$5M
\item Assume the trial identifies the issue before full investment
\item \textbf{Illustrative net avoided cost: HK\$3M}
\end{itemize}\end{block}\end{column}
\end{columns}
\textit{Takeaway: A sandbox trial may create value by discovering problems before scale-up. The result depends on whether the trial is informative.}
\vfill{\tiny Illustrative assumptions; not official Sandbox X costs or timelines.}
\end{frame}

% ====================== Slide 10: CONCEPT PRIMER 3 - BVLOS ======================
\begin{frame}[t]
\frametitle{Concept Primer: BVLOS Operations}
BVLOS means operating an unmanned aircraft beyond the remote pilot's direct visual line of sight.
\begin{itemize}
\item BVLOS can extend range and reduce the need for observers, but usually increases operational complexity and risk.
\item The safety case may need to address command and control, navigation, surveillance, airspace conflict, lost-link procedures, weather, emergency response, and human oversight.
\item Under Hong Kong's risk-based SUA regime, higher-risk advanced operations require prior permission from the Civil Aviation Department and may be subject to additional requirements.
\item Sandbox testing can generate evidence for defined operations, but does not create a general right to conduct BVLOS flights anywhere.
\end{itemize}
\end{frame}
% ====================== Slide 12: SESSION 2 TITLE ======================
\begin{frame}[c]
  \begin{center}
    \huge\color{cuhkpurple}\textbf{Session 2}\\
    \vspace{0.5cm}
    \Large Certification, Safety Cases \& OPEX Savings
  \end{center}
\end{frame}

% ====================== Slide 13: CONCEPT PRIMER - CERTIFICATION ======================
\begin{frame}[t]
\frametitle{Concept Primer: Approval, Permission, and Certification}
These terms should not be used interchangeably.
\begin{itemize}
\item \textbf{SUA operating requirements and permission:} Under Cap. 448G, operations are regulated according to aircraft category and operational risk; advanced operations require prior CAD permission.
\item \textbf{Operator and pilot requirements:} Registration, training, assessment, equipment, records, insurance, and operating conditions may apply.
\item \textbf{Aircraft certification:} Passenger-carrying or larger unconventional aircraft may require a substantially different airworthiness and operational framework.
\item \textbf{Infrastructure approval:} Landing sites and supporting facilities may also involve planning, building, fire, land, environmental, and other requirements.
\end{itemize}
\textbf{Managerial lesson:} Build a regulatory map for the specific aircraft, mission, route, and business model.
\end{frame}

% ====================== Slide 14: APPLIED TOOL - CERTIFICATION ROI ======================
\begin{frame}[t]
\frametitle{Applied Tool: Compliance Investment Appraisal}
\begin{block}{A Simple Expected-Value Screen}
\[
\text{Expected Compliance Value}
=
\Pr(\text{approval})\times\text{Project Value if Approved}
-\text{Compliance Cost}
-\text{Expected Delay Cost}
\]
\end{block}
\begin{itemize}
\item Compliance expenditure may be necessary for market access, but it does not automatically create additional revenue.
\item Project value depends on demand, capacity, competition, operating conditions, and the duration of any advantage.
\item Delay cost should be measured as foregone contribution or project value, not automatically as lost revenue.
\item Use staged decision gates where evidence can change the decision to continue.
\end{itemize}
\end{frame}

% ====================== Slide 15: BACK-OF-THE-ENVELOPE - CERTIFICATION ======================
\begin{frame}[t,shrink=11]
\frametitle{Back-of-the-Envelope: Should the Operator Continue?}
\textbf{Illustrative assumptions:} Compliance programme costs HK\$15M; project value if approved = HK\$30M after all operating costs.
\begin{columns}[T]
\begin{column}{0.48\textwidth}\begin{block}{70\% Approval Probability}
\begin{itemize}
\item Expected approved value: $0.70\times30=\text{HK\$21M}$
\item Less compliance cost: HK\$15M
\item \textbf{Expected value before delay cost: HK\$6M}
\end{itemize}\end{block}\end{column}
\begin{column}{0.48\textwidth}\begin{block}{40\% Approval Probability}
\begin{itemize}
\item Expected approved value: $0.40\times30=\text{HK\$12M}$
\item Less compliance cost: HK\$15M
\item \textbf{Expected value before delay cost: $-$HK\$3M}
\end{itemize}\end{block}\end{column}
\end{columns}
\textit{Takeaway: Compliance investment should be assessed together with approval probability and project economics, not presumed monopoly revenue.}
\vfill{\tiny Illustrative assumptions; not official costs or approval probabilities.}
\end{frame}

% ====================== Slide 16: CONCEPT PRIMER - SAFETY CASE ======================
\begin{frame}[t]
\frametitle{Concept Primer: The Safety Case}
A safety case is a structured argument, supported by evidence, that the risks of a defined operation are understood and adequately controlled.
\begin{itemize}
\item It identifies hazards, consequences, preventive controls, mitigations, responsibilities, and evidence.
\item It should match the actual aircraft, route, environment, operating concept, and emergency procedures.
\item It does not claim to identify every possible failure or prove that risk is zero.
\item It must remain current as operations, technology, routes, and evidence change.
\end{itemize}
\textbf{Business relevance:} A credible safety case can support approval, operational discipline, partner confidence, and insurance discussions.
\end{frame}

% ====================== Slide 17: APPLIED TOOL - INSURANCE DISCOUNT ======================
\begin{frame}[t]
\frametitle{Applied Tool: Risk-Cost Framework}
\begin{block}{Expected Loss}
\[
\text{Expected Loss}=\sum_j \Pr(\text{event }j)\times\text{Financial Consequence}_j
\]
\end{block}
\begin{block}{Total Risk Cost}
\[
\text{Risk Cost}=\text{Prevention and Mitigation Cost}+\text{Insurance Cost}+\text{Retained Expected Loss}
\]
\end{block}
\begin{itemize}
\item Safety investment may reduce event probability or consequence.
\item Insurance premiums also depend on limits, deductibles, exclusions, exposure, claims history, market conditions, and insurer assessment.
\item A stronger safety case may support insurability, but a specific premium reduction is not guaranteed.
\end{itemize}
\end{frame}

% ====================== Slide 18: BACK-OF-THE-ENVELOPE - SAFETY CASE ROI ======================
\begin{frame}[t,shrink=11]
\frametitle{Back-of-the-Envelope: Comparing Risk Controls}
\textbf{Illustrative assumptions:} A proposed control costs HK\$3M and reduces annual expected loss from HK\$2M to HK\$0.5M over a three-year planning period.
\begin{columns}[T]
\begin{column}{0.48\textwidth}\begin{block}{Without the Control}
\begin{itemize}
\item Annual expected loss: HK\$2M
\item Three-year expected loss: HK\$6M
\end{itemize}\end{block}\end{column}
\begin{column}{0.48\textwidth}\begin{block}{With the Control}
\begin{itemize}
\item Control cost: HK\$3M
\item Three-year expected loss: HK\$1.5M
\item Total: HK\$4.5M
\item \textbf{Illustrative expected saving: HK\$1.5M}
\end{itemize}\end{block}\end{column}
\end{columns}
\textit{Takeaway: Safety investment can create economic value when it reduces expected loss or enables operation, but assumptions and non-financial safety obligations remain critical.}
\vfill{\tiny Illustrative risk calculation; not an insurance quotation or acceptable-risk threshold.}
\end{frame}

% ====================== Slide 19: SESSION 3 TITLE ======================
\begin{frame}[c]
  \begin{center}
    \huge\color{cuhkpurple}\textbf{Session 3}\\
    \vspace{0.5cm}
    \Large U-Space, Transaction Costs \& Global Benchmarks
  \end{center}
\end{frame}

% ====================== Slide 20: CONCEPT PRIMER - U-SPACE \& TRANSACTION COSTS ======================
\begin{frame}[t]
\frametitle{Concept Primer: U-space, UTM, and Coordination Cost}
U-space is the European regulatory framework for digital services supporting unmanned-aircraft operations in designated airspace.
\begin{itemize}
\item Mandatory services include flight authorization, geo-awareness, network identification, and traffic information.
\item UTM is a broader term used for systems and services supporting unmanned traffic management in other contexts.
\item Digital services can standardize information exchange and support higher volumes, but they do not eliminate human oversight, certification, cybersecurity, infrastructure, or service-provider costs.
\item Transaction costs include planning, authorization, coordination, information exchange, monitoring, and handling exceptions.
\end{itemize}
\end{frame}

% ====================== Slide 21: APPLIED TOOL - COORDINATION COST ======================
\begin{frame}[t]
\frametitle{Applied Tool: Coordination Cost per Operation}
\begin{block}{A Simple Cost Model}
\[
\text{Daily Coordination Cost}
=
\text{Fixed System Cost}
+q\times\text{Variable Digital Cost}
+h\times\text{Human Oversight Cost}
\]
\end{block}
\begin{itemize}
\item $q$ is the number of operations and $h$ is the required human-oversight effort.
\item Digitalization may reduce repetitive manual work while adding software, communications, cybersecurity, certification, and provider fees.
\item Exceptions, disruptions, and emergency situations may still require substantial human intervention.
\item Scale is viable only if total system capacity and safety performance also increase.
\end{itemize}
\end{frame}

% ====================== Slide 22: BACK-OF-THE-ENVELOPE - U-SPACE ======================
\begin{frame}[t,shrink=11]
\frametitle{Back-of-the-Envelope: Manual and Digital Coordination}
\textbf{Assumptions:} 1,000 operations/day and 30 operating days/month.
\begin{columns}[T]
\begin{column}{0.48\textwidth}\begin{block}{Mostly Manual Process}
\begin{itemize}
\item 100 staff-hours/day at HK\$300/hour
\item Daily cost: HK\$30,000
\item \textbf{Monthly cost: HK\$900,000}
\end{itemize}\end{block}\end{column}
\begin{column}{0.48\textwidth}\begin{block}{Digital System with Oversight}
\begin{itemize}
\item System cost: HK\$8,000/day
\item Human oversight: HK\$7,000/day
\item Daily cost: HK\$15,000
\item \textbf{Monthly cost: HK\$450,000}
\end{itemize}\end{block}\end{column}
\end{columns}
\textit{Takeaway: Digital coordination may lower unit cost, but savings depend on system cost, oversight, reliability, integration, and scale.}
\vfill{\tiny Illustrative assumptions; not a U-space or Hong Kong UTM cost estimate.}
\end{frame}

% ====================== Slide 23: CONCEPT PRIMER - GLOBAL BENCHMARKS ======================
\begin{frame}[t]
\frametitle{Concept Primer: Comparing Regulatory Systems}
International comparison should focus on functions rather than unsupported rankings of speed.
\begin{itemize}
\item \textbf{Hong Kong:} A risk-based SUA framework, prior permission for advanced operations, and controlled Sandbox X trials for more complex applications.
\item \textbf{European Union:} Harmonized UAS rules and a U-space framework for designated airspace and certified service providers.
\item \textbf{United States:} Certification, operational approvals, airspace integration, infrastructure, and implementation planning under FAA and interagency programmes.
\item \textbf{Mainland China and the GBA:} A distinct national and local regulatory environment requiring route-specific and cross-boundary coordination.
\end{itemize}
\textbf{Managerial lesson:} Compare requirements, evidence, scope, transferability, cost, and predictability, not only nominal approval time.
\end{frame}
% ====================== Slide 24: APPLIED TOOL - REGULATORY AGILITY INDEX ======================
\begin{frame}[t]
\frametitle{Applied Tool: Regulatory-Pathway Scorecard}
\begin{center}\small
\begin{tabular}{l|c}
\textbf{Dimension} & \textbf{Managerial question}\\
\hline
Clarity & Are requirements and decision points understandable?\\
Proportionality & Are controls matched to operational risk?\\
Evidence pathway & Can testing generate useful evidence?\\
Predictability & Are process, timing, and responsibilities credible?\\
Scalability & Can the approval support expansion or repetition?\\
Interoperability & Can evidence or approvals support other markets?\\
\end{tabular}
\end{center}
\vspace{0.25cm}
\textit{A single ratio of approval time to an invented ``stringency score'' would not provide a reliable international comparison.}
\end{frame}

% ====================== Slide 25: BACK-OF-THE-ENVELOPE - GLOBAL BENCHMARKS ======================
\begin{frame}[t]
\frametitle{Back-of-the-Envelope: Weighted Pathway Comparison}
\textbf{Illustrative exercise:} Score each pathway from 1 to 5 using evidence available to the project team.
\begin{center}\small
\begin{tabular}{l|c|c}
\textbf{Dimension} & \textbf{Weight} & \textbf{What the team documents}\\
\hline
Clarity and predictability & 25\% & Published rules and milestones\\
Trial and evidence pathway & 20\% & Scope and usefulness of testing\\
Time and internal effort & 20\% & Scenario range, not a fixed claim\\
Scalability & 20\% & Route, fleet, and repeatability limits\\
Transferability & 15\% & Recognition and re-use of evidence\\
\end{tabular}
\end{center}
\[
\text{Weighted Score}=\sum_k w_k s_k
\]
\textit{Takeaway: The scorecard makes assumptions visible. It should support judgment rather than claim that one jurisdiction is universally fastest or best.}
\end{frame}

% ====================== Slide 26: CONCEPT PRIMER - POLICY HARMONISATION ======================
\begin{frame}[t]
\frametitle{Concept Primer: Coordination and Harmonisation}
Cross-boundary operations may require coordination across multiple legal and operational systems.
\begin{itemize}
\item Harmonisation means aligning selected rules, standards, data formats, procedures, or evidence requirements.
\item Mutual recognition means one authority accepts specified approvals or evidence from another jurisdiction.
\item Neither concept necessarily creates one licence valid everywhere.
\item Cross-boundary operations may also require customs, immigration, quarantine, spectrum, data, security, liability, and emergency-response arrangements.
\end{itemize}
\textbf{Managerial lesson:} Identify exactly which requirement must be aligned, recognized, duplicated, or separately approved.
\end{frame}


% ====================== Slide 18: EXERCISE INSTRUCTIONS ======================
\begin{frame}[t]
\frametitle{In-Class Exercise: Regulatory Simulation}
\textbf{Scenario:} An operator proposes a drone-delivery route near a residential area in Sha Tin.
\begin{itemize}
\item \textbf{Format:} Operator, regulator, public-agency customer, and community representatives
\item \textbf{Task:}
\begin{enumerate}
\item define the service value and affected stakeholders;
\item identify the main operational and community risks;
\item propose initial route, time, payload, weather, monitoring, and contingency conditions;
\item specify the evidence required before expanding the operation; and
\item define a pause, review, and incident-response rule.
\end{enumerate}
\item \textbf{Timing:} 15 minutes plus 5-minute discussion.
\end{itemize}
\end{frame}

% ====================== Slide 19: EXERCISE TEMPLATE ======================
\begin{frame}[t]
\frametitle{Exercise Template: Evidence-Based Roadmap}
\begin{center}\small
\begin{tabular}{p{0.18\textwidth}|p{0.23\textwidth}|p{0.25\textwidth}|p{0.23\textwidth}}
\textbf{Phase} & \textbf{Operating scope} & \textbf{Evidence collected} & \textbf{Expansion gate}\\
\hline
1. Controlled test & Limited route, time, and payload & Communications, navigation, weather, procedures & Safety and reliability thresholds met\\
2. Extended trial & More operating conditions or users & Repeated performance and stakeholder feedback & Approved review and updated controls\\
3. Scaled operation & Wider route, frequency, or fleet & Ongoing monitoring and incident reporting & Continued compliance and periodic review\\
\end{tabular}
\end{center}
\vspace{0.25cm}
\textit{Data sharing should be necessary, proportionate, secure, and governed by clear rules. It should not be offered as an informal exchange for faster approval.}
\end{frame}

% ====================== Slide 20: KEY TAKEAWAYS ======================
\begin{frame}[t]
\frametitle{Key Takeaways}
\begin{itemize}
\item Hong Kong regulates small unmanned aircraft through a risk-based framework; advanced operations require prior permission and may face additional conditions.
\item Sandbox X supports controlled testing and evidence collection, but is not a waiver of safety rules or automatic commercial approval.
\item Compliance capability can create value through readiness, learning, reliability, and credibility, not guaranteed monopoly protection.
\item A safety case is an evidence-based argument for a defined operation and may support approval and insurability without guaranteeing lower premiums.
\item Digital traffic-management services may reduce coordination cost and support scale, but still require oversight, cybersecurity, infrastructure, and governance.
\item A strong roadmap links each expansion in operating privilege to specified evidence, review, and risk controls.
\end{itemize}
\end{frame}

% ====================== Slide 21: ASSIGNMENT ======================
\begin{frame}[t]
\frametitle{Week 7 Assignment Briefing}
\textbf{Regulatory and Compliance Roadmap Memo (max 1,200 words)}
\begin{itemize}
\item \textbf{Due:} Beginning of Week 8 (group assignment)
\item \textbf{Task:} Develop a regulatory strategy for the group's Hong Kong LAE use case.
\item \textbf{Requirements:}
\begin{enumerate}
\item identify the aircraft, operation, route, stakeholders, and applicable regulatory pathway;
\item distinguish registration, operating permission, aircraft approval, and site requirements;
\item outline the major hazards and evidence required for a safety case;
\item propose a three-phase test-to-scale roadmap with measurable expansion gates;
\item appraise the expected value and cost of early testing using transparent assumptions; and
\item propose a necessary and proportionate data-governance plan for regulatory reporting.
\end{enumerate}
\item \textbf{Grading focus:} Regulatory accuracy (30\%), economic reasoning (30\%), roadmap and risk controls (30\%), professionalism (10\%).
\end{itemize}
\end{frame}

% ====================== Slide 22: THANK YOU ======================
\begin{frame}[t]
  \frametitle{Thank You + Q\&A}
  
  \begin{center}
      \textbf{Next Week Preview:} \\
      \Large Airspace Access and Integration
  \end{center}
  
  \vspace{0.5cm}
  \textbf{Pre-Class Readings Reminder (for Week 8):}
  \begin{itemize}
      \item \textit{CAD AC-013 (BVLOS) and Safety Requirements Document}
      \item \textit{EASA Easy Access Rules for U-space}
  \end{itemize}
  
  \vspace{0.5cm}
  \begin{center}
      \textbf{Questions?} \\
      \small (Final 15 minutes reserved for extended Q\&A)
  \end{center}
\end{frame}

\end{document}